\documentclass[../../../include/open-logic-section]{subfiles}

\begin{document}
\olfileid[pt]{sth}{ord-arithmetic}{intro}

\olsection{Introdução}

No \olref[ordinals][]{chap}, desenvolvemos uma teoria dos números ordinais.
Vimos no \olref[spine][]{chap} que podemos pensar nos ordinais como uma
espinha dorsal em torno da qual se constrói o restante da hierarquia. Mas
esse não é o único papel dos ordinais. Há também a tarefa de
efetuar a aritmética ordinal.

Já aludimos a isto, em
\olref[ordinals][idea]{sec}, quando falamos de $\omega$,
$\omega+1$ e $\omega+\omega$. Na altura falamos informalmente; chegou
a altura de o expor com rigor. Devemos contudo mencionar
que não há muita filosofia neste capítulo; apenas desenvolvimentos
técnicos, juntamente com uma observação (ligeiramente) interessante de que
podemos fazer a mesma coisa de duas maneiras diferentes.

\end{document}
